\scanpage{14}

$K$ noyau algébrique sur $k$ corps du Thm.

On veut prouver $H^2(k,K)\simeq H^2(k,K)'\simeq\text{un élément}$ [résulte formellement de \uuline{$H^2(k,K)'\ne\varnothing$}, compte tenu que $H^2(k,\unsure{Z|\mathfrak{z}}(K))=0$ !] \deleted{\unsure{Raisonnons sur} \unsure{dimension} de $K$, \unsure{trivial si elle est}}

a) $\dim K=0$. On raisonne par récurrence sur le degré de $K$. Trivial s'il est égal à $1$. S'il est $>1$, par dévissage \inserted{($2^{\text{ième}}$ principe de réduction)}, on est ramené aux cas suivants

(i) $K$ étale et \unsure{on|en} aux deux cas
\[
\left|\begin{aligned}
&K\ \deleted{\text{étale}}\ \text{comm. (trivial)}\\
&[K,K]=K\quad\text{en particulier }K\ \uline{\text{sans centre}}.
\end{aligned}\right.
\]

(ii) $K$ radiciel et $=$ par dévissage
\[
\left\{\begin{aligned}
&K\ \text{radiciel de hauteur }1\\
&\text{et }=\ [K,K]=1\quad\text{en particulier }K\ \text{sans centre}.
\end{aligned}\right.
\]

Le cas (i) correspond à un hom de $\pi$ (groupe de Galois top. de $\bar k/k$] dans $\operatorname{Aut}\operatorname{ext}(A)$, $A$ groupe fini constant. On sait qu'il se remonte à un hom $\pi\unsure{-|\to}\operatorname{Aut}(A)$, [\unsure{Schur|Schreier}+Serre] donc on gagne \unsure{...}

Dans le cas (ii), on considère les sections $\unsure{\varphi|\phi}$ de $\unsure{\sigma|\mathscr{T}}(K)$ correspondant aux ``ss-groupes de Cartan \inserted{infinitésimaux}'' \deleted{qui \unsure{sont leurs}} qui satisfont $\underline{\unsure{\operatorname{Nor}|\operatorname{Na}}}(\varphi)=\varphi$. Par le principe de réduction $1$, on est

% archive sheet number: 13

